\chapter{História e Geografia do Estado do Maranhão}

\section{Como a banca transforma conteúdo regional em questão seletiva}

História e Geografia do Maranhão não devem ser estudadas como coleção de datas e nomes. Em prova de alto nível, o examinador costuma explorar \textbf{causa, consequência, sequência temporal, localização, relação entre meio físico e economia e comparação entre processos}. A técnica de estudo é sempre responder quatro perguntas: \emph{o que ocorreu? por que ocorreu? em que contexto? que efeito produziu?}

\begin{teoria}[linha mestra]
O conteúdo do edital pode ser organizado em três eixos integrados:
\begin{enumerate}
\item \textbf{formação histórica}: França Equinocial, disputa colonial, Estado do Maranhão, revoltas, Independência e conflitos sociais;
\item \textbf{território e ambiente}: posição geográfica, relevo, clima, hidrografia, vegetação e unidades de conservação;
\item \textbf{ocupação e economia}: população, urbanização, agricultura, pecuária, extrativismo, indústria, logística, portos e cultura.
\end{enumerate}
\end{teoria}

\section{França Equinocial e fundação de São Luís}

No início do século XVII, franceses buscaram estabelecer uma colônia no norte da América portuguesa. A expedição liderada por Daniel de La Touche, senhor de La Ravardière, instalou-se na ilha do Maranhão em 1612. O projeto ficou conhecido como \termo{França Equinocial}. A ocupação francesa não deve ser confundida com a França Antártica, experiência anterior na baía de Guanabara.

A fundação de São Luís está associada a essa presença francesa em 1612. O domínio francês, porém, foi breve. A reação luso-portuguesa resultou em confrontos, entre os quais a \termo{Batalha de Guaxenduba}, em 1614, e na consolidação portuguesa da região.

\begin{cebraspe}
Pegadinha clássica: trocar França Equinocial por França Antártica, ou deslocar a Batalha de Guaxenduba para o contexto da invasão holandesa. Guaxenduba pertence ao conflito luso-francês do início do século XVII.
\end{cebraspe}

\section{Invasão holandesa e expulsão}

A presença neerlandesa no Maranhão ocorreu no contexto mais amplo das disputas atlânticas do século XVII e da expansão holandesa sobre áreas coloniais portuguesas. São Luís foi ocupada em 1641. A resistência local e luso-portuguesa culminou na expulsão dos holandeses em 1644.

O ponto importante para a prova é perceber que \textbf{franceses e holandeses correspondem a ciclos distintos}: os franceses estão ligados à fundação de São Luís e à França Equinocial; os holandeses ocupam a região décadas depois, em outro contexto geopolítico.

\section{Estado do Maranhão e Grão-Pará}

A administração colonial do norte do território português foi organizada de modo distinto do Estado do Brasil em diferentes períodos. A criação de unidades administrativas próprias buscava responder à distância, às dificuldades de comunicação e à importância estratégica da Amazônia e do litoral setentrional.

Esse arranjo ajuda a explicar peculiaridades econômicas e políticas da região, bem como conflitos entre colonos, Coroa, missionários e companhias monopolistas.

\subsection{Revolta de Beckman}

A Revolta de Beckman ocorreu em 1684, em São Luís. Entre as causas estavam insatisfação de colonos com o monopólio e o desempenho da Companhia de Comércio do Maranhão, problemas de abastecimento e conflitos relacionados à mão de obra indígena e à atuação jesuítica.

O movimento contestou mecanismos de administração econômica colonial, mas não foi uma revolta de independência nacional. Essa distinção é importante: a crítica incidia sobre práticas e agentes do sistema colonial, sem formular projeto nacional soberano nos moldes do século XIX.

\section{Independência e adesão do Maranhão}

A Independência proclamada em 1822 não produziu adesão imediata e uniforme em todas as províncias. No Maranhão, havia fortes vínculos políticos e econômicos com Portugal. A incorporação da província ao novo Império brasileiro ocorreu em 1823, após conflitos militares e pressão das forças favoráveis à Independência.

A \termo{Batalha do Jenipapo}, travada em 1823 no Piauí, integra esse processo regional de consolidação da Independência no norte do Brasil. Embora ocorrida no território piauiense, possui conexão histórica com a adesão das províncias vizinhas, inclusive o Maranhão.

\begin{atencao}
Não trate a Independência como evento instantâneo em 7 de setembro de 1822. Para o Maranhão, o examinador pode explorar precisamente a diferença entre proclamação nacional e adesão provincial.
\end{atencao}

\section{Balaiada}

A Balaiada, entre 1838 e 1841, foi uma revolta social ocorrida principalmente no Maranhão e no Piauí durante o Período Regencial. Envolveu sertanejos, vaqueiros, artesãos, escravizados e outros grupos populares, em cenário de pobreza, tensões políticas locais e disputa entre elites.

O nome está associado a Manuel Francisco dos Anjos Ferreira, o Balaio. Entre os personagens do movimento também se destacam Raimundo Gomes e Cosme Bento. A repressão foi conduzida por forças imperiais sob comando de Luís Alves de Lima e Silva, futuro Duque de Caxias.

A prova pode exigir diferenciação entre causas sociais profundas e disputas políticas imediatas. Reduzir a Balaiada a uma simples revolta partidária elimina a participação popular e a estrutura socioeconômica do conflito.

\section{República, 1930, Vitorinismo e Greve de 1951}

A transição para a República também repercutiu nas estruturas políticas locais. Ao longo da Primeira República, elites estaduais e redes oligárquicas tiveram papel central na organização do poder.

A Revolução de 1930 rompeu o arranjo político da República Velha e alterou as relações entre governo federal e elites estaduais. No Maranhão, o período posterior foi marcado por reconfigurações políticas, interventorias e novas disputas de poder.

O \termo{Vitorinismo} está associado à liderança política de Vitorino Freire e à forte influência de seu grupo sobre a política maranhense, sobretudo entre as décadas de 1940 e 1960. A \termo{Greve de 1951} ocorreu em ambiente de contestação político-eleitoral e mobilização popular, sendo marco relevante das tensões daquele período.

Na segunda metade do século XX, a política maranhense passou por mudanças de grupos dominantes, modernização parcial da infraestrutura, crescimento urbano, expansão de grandes projetos e persistência de desigualdades socioeconômicas.

\section{Localização e organização territorial}

O Maranhão situa-se na Região Nordeste, em posição de transição entre diferentes domínios naturais e econômicos. Limita-se com o oceano Atlântico ao norte, Piauí a leste, Tocantins ao sul/sudoeste e Pará a oeste.

Essa posição ajuda a explicar sua diversidade ambiental: áreas de influência amazônica no oeste, cerrado no centro-sul, extensas formações de babaçuais e ambientes costeiros com manguezais, dunas e estuários.

\section{Clima, relevo, geologia e recursos minerais}

O clima maranhense varia espacialmente. O norte e o oeste são, em geral, mais úmidos, enquanto o sul e o sudeste apresentam sazonalidade mais marcada, com estação seca mais definida. A pluviosidade e a temperatura condicionam vegetação, regime dos rios e atividades agropecuárias.

O relevo combina planícies litorâneas e fluviais, baixadas e superfícies de planalto. Em questões, associe as formas do relevo a processos geomorfológicos, drenagem, ocupação e uso econômico, em vez de decorar listas isoladas.

A geologia condiciona ocorrência de recursos minerais e materiais utilizados economicamente. O edital exige compreender a relação entre estrutura geológica, recursos minerais e atividade extrativa, sem restringir o estudo a um único produto.

\section{Hidrografia}

O Maranhão possui rede hidrográfica extensa. O edital separa \textbf{bacias limítrofes} e \textbf{bacias genuinamente maranhenses}.

\begin{table}[ht]
\centering\small
\begin{tabularx}{\textwidth}{l X}
\toprule
Grupo & Exemplos e leitura de prova \\
\midrule
Bacias limítrofes & Parnaíba, Gurupi e sistema Tocantins-Araguaia; cursos associados a limites ou integração interestadual. \\
Bacias maranhenses & Mearim, Itapecuru, Munim, Pericumã e outras drenagens internas/estaduais; o Mearim e o Itapecuru são especialmente relevantes para ocupação e abastecimento. \\
\bottomrule
\end{tabularx}
\end{table}

A hidrografia deve ser relacionada ao abastecimento urbano, agricultura, transporte local, ecossistemas e risco de cheias. A Baixada Maranhense, por exemplo, possui forte dinâmica sazonal de inundação.

\section{Vegetação e unidades de conservação}

O estado apresenta zonas de transição entre Floresta Amazônica, Cerrado e formações de cocais, além de manguezais e vegetação costeira. A \termo{Mata dos Cocais} é formação de transição marcada especialmente por palmeiras como babaçu e carnaúba.

Entre as áreas protegidas, destacam-se unidades de proteção integral e de uso sustentável. O edital cita APAs e parques nacionais. Para prova, diferencie a lógica das categorias: parque nacional é unidade de proteção integral; APA é unidade de uso sustentável com ocupação e atividades condicionadas a regras de proteção.

Os Lençóis Maranhenses constituem uma das paisagens mais conhecidas do estado, formada por campos de dunas e lagoas sazonais sob dinâmica climática e costeira específica. A proteção ambiental deve ser analisada conjuntamente com turismo, comunidades locais e conservação.

\section{População, urbanização e movimentos populacionais}

A população maranhense distribui-se de maneira desigual. São Luís concentra funções administrativas, portuárias, industriais, comerciais e de serviços. Imperatriz exerce centralidade regional no sudoeste, articulando fluxos com Tocantins e Pará.

Urbanização não significa distribuição homogênea de infraestrutura. Questões podem relacionar crescimento urbano a migração interna, expansão periférica, desigualdade de acesso a serviços e integração regional.

\section{Agricultura, pecuária e extrativismo}

A agricultura maranhense combina produção familiar, culturas tradicionais e agricultura empresarial. Arroz, mandioca, milho e feijão possuem importância histórica, enquanto soja e outros grãos ganharam grande peso no sul e leste do estado, associados à expansão do agronegócio no MATOPIBA.

A pecuária bovina é relevante em diversas regiões. O extrativismo vegetal, especialmente do babaçu, possui importância econômica, social e cultural. O extrativismo mineral e atividades vinculadas a grandes cadeias logísticas também integram a economia estadual.

\section{Indústria, logística e setor terciário}

O Maranhão possui localização estratégica para corredores de exportação. O Complexo Portuário de São Luís e o Porto do Itaqui articulam cargas minerais, agrícolas, combustíveis e outros produtos. A Estrada de Ferro Carajás conecta a região de Carajás ao litoral maranhense, influenciando fortemente a logística e a atividade industrial-portuária.

O setor terciário concentra parcela expressiva das atividades urbanas: comércio, serviços, telecomunicações, administração pública e transportes. A malha rodoviária e ferroviária, os portos e aeroportos condicionam integração territorial e competitividade.

\section{Cultura maranhense}

A cultura do Maranhão resulta de matrizes indígenas, africanas e europeias e de processos históricos próprios. Entre manifestações de alta relevância estão o \termo{bumba-meu-boi}, o \termo{tambor de crioula}, festas religiosas, culinária, artesanato, azulejaria de São Luís e a forte presença do reggae na identidade urbana da capital.

Em prova, cultura pode aparecer integrada a território e história: manifestações culturais não são elementos isolados, mas expressão das relações sociais e das múltiplas matrizes de formação do estado.

\section{Mapa mental funcional}

\mapafuncional{Maranhão — História e Geografia}
{Resolver questão regional = localizar o período/território + identificar causa + relacionar consequência}
{1612 França Equinocial $\rightarrow$ 1614 Guaxenduba.\\1641--1644 ocupação holandesa.\\1684 Beckman.\\1823 adesão à Independência.\\1838--1841 Balaiada.\\Século XX: 1930, Vitorinismo, 1951.}
{Fronteiras: Atlântico/Piauí/Tocantins/Pará.\\Ambientes: floresta, cerrado, cocais, mangues.\\Rios: Parnaíba/Gurupi/Tocantins-Araguaia vs. Mearim/Itapecuru/Munim.\\Economia: agropecuária + extrativismo + logística portuária.}
{França Equinocial $\neq$ França Antártica.\\Guaxenduba $\neq$ invasão holandesa.\\Beckman $\neq$ independência.\\Balaiada $\neq$ revolta apenas de elites.\\APA $\neq$ parque nacional.}
{Qual período? Qual agente? Qual conflito? Qual espaço físico?\\Há relação entre clima/relevo/vegetação?\\O rio é limítrofe ou maranhense?\\A atividade econômica depende de qual corredor logístico?}

\begin{resumo}
\textbf{Revisão de 60 segundos:} França Equinocial e fundação de São Luís; Guaxenduba; holandeses; Beckman; adesão em 1823; Balaiada; 1930/Vitorinismo/1951; transição Amazônia--Cerrado--Cocais; Parnaíba/Gurupi/Tocantins-Araguaia; Mearim/Itapecuru; babaçu; agronegócio; Itaqui e corredores ferroviários; cultura do bumba-meu-boi e tambor de crioula.
\end{resumo}

\section{Banco de questões — História e Geografia do Maranhão}

\subsection*{N1 — Fundamentos em nível de concurso}
\begin{questao}{\nivel{N1} 08.01 — Cebraspe (C/E)}A França Equinocial, instalada em 1612, relaciona-se diretamente à fundação de São Luís e não deve ser confundida com a França Antártica.\end{questao}
\begin{questao}{\nivel{N1} 08.02 — Cebraspe (C/E)}A Batalha de Guaxenduba ocorreu no contexto da reação luso-portuguesa à presença francesa no Maranhão.\end{questao}
\begin{questao}{\nivel{N1} 08.03 — FGV}A ocupação holandesa de São Luís ocorreu:\begin{enumerate}[label=\Alph*)]\item antes da França Equinocial.\item no século XVII, posteriormente ao conflito luso-francês.\item durante a Balaiada.\item após a Revolução de 1930.\item no Segundo Reinado.\end{enumerate}\end{questao}
\begin{questao}{\nivel{N1} 08.04 — Cebraspe (C/E)}A Revolta de Beckman esteve associada à insatisfação com mecanismos econômicos coloniais, incluindo a atuação da Companhia de Comércio do Maranhão.\end{questao}
\begin{questao}{\nivel{N1} 08.05 — Cebraspe (C/E)}A adesão do Maranhão à Independência ocorreu de modo imediato em 7 de setembro de 1822, sem resistência ou conflito regional posterior.\end{questao}
\begin{questao}{\nivel{N1} 08.06 — FCC}A Balaiada caracterizou-se por:\begin{enumerate}[label=\Alph*)]\item movimento exclusivamente militar sem participação popular.\item revolta social do Período Regencial com forte participação popular.\item invasão estrangeira.\item rebelião republicana de 1889.\item movimento exclusivamente urbano de São Luís em 1951.\end{enumerate}\end{questao}
\begin{questao}{\nivel{N1} 08.07 — Cebraspe (C/E)}O Maranhão limita-se com Piauí, Tocantins e Pará, além de possuir litoral voltado para o oceano Atlântico.\end{questao}
\begin{questao}{\nivel{N1} 08.08 — Cebraspe (C/E)}Parnaíba, Gurupi e Tocantins-Araguaia aparecem no edital como bacias limítrofes, enquanto Mearim e Itapecuru são exemplos relevantes de bacias maranhenses.\end{questao}
\begin{questao}{\nivel{N1} 08.09 — Vunesp}No sistema de unidades de conservação, uma APA é, em regra:\begin{enumerate}[label=\Alph*)]\item unidade de proteção integral necessariamente sem ocupação humana.\item unidade de uso sustentável.\item área sem qualquer regime de proteção.\item sinônimo de parque nacional.\item área destinada exclusivamente à mineração.\end{enumerate}\end{questao}
\begin{questao}{\nivel{N1} 08.10 — Cebraspe (C/E)}A Mata dos Cocais é formação de transição na qual se destacam palmeiras como o babaçu.\end{questao}

\subsection*{N2 — Aplicação}
\begin{questao}{\nivel{N2} 08.11 — FGV}Uma questão associa Daniel de La Touche, fundação de São Luís e França Equinocial. A relação é:\begin{enumerate}[label=\Alph*)]\item historicamente coerente.\item incorreta, pois La Touche liderou a Balaiada.\item incorreta, pois São Luís foi fundada por holandeses.\item incorreta, pois França Equinocial ocorreu no século XIX.\item incorreta, pois o episódio se deu no Piauí.\end{enumerate}\end{questao}
\begin{questao}{\nivel{N2} 08.12 — Cebraspe (C/E)}A Batalha do Jenipapo, embora travada no Piauí, integra o contexto regional de consolidação da Independência no norte do Brasil e se relaciona à adesão do Maranhão ao Império.\end{questao}
\begin{questao}{\nivel{N2} 08.13 — Cebraspe (C/E)}Reduzir a Revolta de Beckman a movimento de independência nacional seria historicamente inadequado.\end{questao}
\begin{questao}{\nivel{N2} 08.14 — FCC}O Vitorinismo está associado principalmente:\begin{enumerate}[label=\Alph*)]\item à presença francesa no século XVII.\item à influência política de Vitorino Freire no Maranhão do século XX.\item à Guerra do Paraguai.\item à expansão holandesa.\item à criação do Estado do Grão-Pará no século XVII.\end{enumerate}\end{questao}
\begin{questao}{\nivel{N2} 08.15 — Cebraspe (C/E)}A Greve de 1951 deve ser compreendida no contexto de tensões políticas e mobilização popular do Maranhão do pós-guerra.\end{questao}
\begin{questao}{\nivel{N2} 08.16 — FGV}A diversidade ambiental do Maranhão é melhor explicada por sua posição:\begin{enumerate}[label=\Alph*)]\item exclusivamente amazônica.\item exclusivamente sertaneja.\item de transição entre domínios amazônicos, cerrados, cocais e ambientes costeiros.\item inteiramente inserida na Caatinga.\item sem contato com o oceano.\end{enumerate}\end{questao}
\begin{questao}{\nivel{N2} 08.17 — Cebraspe (C/E)}A pluviosidade tende a apresentar diferenças espaciais no Maranhão, influenciando regime fluvial, vegetação e uso agropecuário.\end{questao}
\begin{questao}{\nivel{N2} 08.18 — Cebraspe (C/E)}Mearim e Itapecuru possuem relevância para ocupação, abastecimento e atividades econômicas, razão pela qual sua leitura não deve se limitar à memorização do nome da bacia.\end{questao}
\begin{questao}{\nivel{N2} 08.19 — FGV}A Baixada Maranhense relaciona-se de modo marcante a:\begin{enumerate}[label=\Alph*)]\item dinâmica sazonal de inundação.\item relevo exclusivamente montanhoso.\item ausência de corpos d'água.\item clima desértico.\item mineração de carvão como única atividade.\end{enumerate}\end{questao}
\begin{questao}{\nivel{N2} 08.20 — Cebraspe (C/E)}A expansão da soja no sul e no leste do Maranhão integra a dinâmica recente do agronegócio associada ao MATOPIBA.\end{questao}
\begin{questao}{\nivel{N2} 08.21 — FCC}O Porto do Itaqui e os corredores ferroviários são relevantes porque:\begin{enumerate}[label=\Alph*)]\item reduzem a integração do estado com mercados externos.\item articulam fluxos de exportação, combustíveis, minérios e produtos agrícolas.\item possuem função apenas turística.\item substituem toda a malha rodoviária.\item se localizam fora do Maranhão.\end{enumerate}\end{questao}
\begin{questao}{\nivel{N2} 08.22 — Cebraspe (C/E)}A concentração de serviços e funções administrativas em São Luís ajuda a explicar sua centralidade urbana no estado.\end{questao}
\begin{questao}{\nivel{N2} 08.23 — Cebraspe (C/E)}Imperatriz exerce importante centralidade regional e articula fluxos com estados vizinhos, especialmente na porção sudoeste do Maranhão.\end{questao}
\begin{questao}{\nivel{N2} 08.24 — FGV}Babaçu, bumba-meu-boi e Porto do Itaqui pertencem, respectivamente, aos eixos:\begin{enumerate}[label=\Alph*)]\item extrativismo vegetal, cultura e logística.\item mineração, pecuária e religião.\item agricultura irrigada, indústria automobilística e pesca.\item cultura, hidrografia e climatologia.\item política, relevo e demografia.\end{enumerate}\end{questao}
\begin{questao}{\nivel{N2} 08.25 — Cebraspe (C/E)}Urbanização e crescimento populacional não implicam, por si, distribuição homogênea de infraestrutura e serviços públicos.\end{questao}

\subsection*{N3 — Nível de prova}
\begin{questao}{\nivel{N3} 08.26 — Cebraspe (C/E)}Associar a Batalha de Guaxenduba à expulsão dos holandeses é incorreto, pois o episódio integra o conflito contra os franceses.\end{questao}
\begin{questao}{\nivel{N3} 08.27 — FGV}A melhor sequência cronológica é:\begin{enumerate}[label=\Alph*)]\item Balaiada -- França Equinocial -- Beckman -- Independência.\item França Equinocial -- Guaxenduba -- ocupação holandesa -- Beckman.\item Beckman -- Guaxenduba -- Balaiada -- França Equinocial.\item Independência -- França Equinocial -- holandeses -- Balaiada.\item 1930 -- Balaiada -- Beckman -- Guaxenduba.\end{enumerate}\end{questao}
\begin{questao}{\nivel{N3} 08.28 — Cebraspe (C/E)}A Balaiada combina conflitos entre elites locais e forte participação de grupos populares, não sendo adequadamente explicada por uma única causa.\end{questao}
\begin{questao}{\nivel{N3} 08.29 — Cebraspe (C/E)}A Revolução de 1930 alterou relações políticas entre poder central e elites estaduais, afetando também a organização política do Maranhão.\end{questao}
\begin{questao}{\nivel{N3} 08.30 — FGV}Considere: I. França Equinocial; II. Revolta de Beckman; III. Balaiada. Os três eventos pertencem, respectivamente, aos contextos:\begin{enumerate}[label=\Alph*)]\item disputa colonial; crise do sistema colonial; Período Regencial.\item República; Império; Colônia.\item Império; República; Colônia.\item Estado Novo; Colônia; República.\item todos ao século XX.\end{enumerate}\end{questao}
\begin{questao}{\nivel{N3} 08.31 — Cebraspe (C/E)}A posição do Maranhão entre Amazônia e Nordeste contribui para mosaico de vegetação e para diferenças internas de precipitação e uso da terra.\end{questao}
\begin{questao}{\nivel{N3} 08.32 — FGV}Se uma área apresenta babaçuais extensos em zona de transição, a formação vegetal mais diretamente associada é:\begin{enumerate}[label=\Alph*)]\item Mata dos Cocais.\item tundra.\item pradaria pampeana.\item araucária.\item pantanal.\end{enumerate}\end{questao}
\begin{questao}{\nivel{N3} 08.33 — Cebraspe (C/E)}Os manguezais do litoral maranhense possuem relação com ambientes estuarinos e costeiros e não podem ser classificados simplesmente como cerrado.\end{questao}
\begin{questao}{\nivel{N3} 08.34 — FCC}A diferença fundamental entre APA e parque nacional reside em que:\begin{enumerate}[label=\Alph*)]\item ambos são idênticos.\item APA é categoria de uso sustentável, enquanto parque nacional integra proteção integral.\item parque nacional é sempre privado.\item APA não possui objetivo ambiental.\item parque nacional admite qualquer exploração mineral sem restrição.\end{enumerate}\end{questao}
\begin{questao}{\nivel{N3} 08.35 — Cebraspe (C/E)}Uma política de uso do solo no Maranhão deve considerar que relevo, drenagem e regime de chuvas influenciam erosão, inundação e aptidão agrícola.\end{questao}
\begin{questao}{\nivel{N3} 08.36 — FGV}A Estrada de Ferro Carajás possui importância para o Maranhão porque:\begin{enumerate}[label=\Alph*)]\item articula corredor mineral e logístico até a região portuária de São Luís.\item conecta exclusivamente cidades do litoral piauiense.\item é ferrovia urbana de passageiros da capital.\item substitui o Porto do Itaqui.\item não possui relação com exportações.\end{enumerate}\end{questao}
\begin{questao}{\nivel{N3} 08.37 — Cebraspe (C/E)}A economia maranhense combina produção familiar, agronegócio de grãos, pecuária, extrativismo, indústria e serviços, de modo que classificá-la por um único setor é reducionista.\end{questao}
\begin{questao}{\nivel{N3} 08.38 — Cebraspe (C/E)}A expansão de grandes corredores logísticos pode gerar crescimento econômico sem eliminar automaticamente desigualdades territoriais e sociais.\end{questao}
\begin{questao}{\nivel{N3} 08.39 — FGV}Em análise regional, São Luís e Imperatriz devem ser vistas como:\begin{enumerate}[label=\Alph*)]\item centros urbanos com funções e áreas de influência distintas.\item municípios sem função regional.\item áreas rurais exclusivamente agrícolas.\item portos marítimos equivalentes.\item cidades sem conexão rodoviária.\end{enumerate}\end{questao}
\begin{questao}{\nivel{N3} 08.40 — Cebraspe (C/E)}Bumba-meu-boi e tambor de crioula expressam a formação cultural plural do Maranhão e não devem ser reduzidos a manifestações de uma única matriz étnica.\end{questao}

\subsection*{N4 — Avançado e integração}
\begin{questao}{\nivel{N4} 08.41 — FGV}Uma questão relaciona clima mais sazonal no sul do Maranhão, expansão de grãos e integração por corredores logísticos. A relação é plausível porque:\begin{enumerate}[label=\Alph*)]\item condições naturais e infraestrutura influenciam localização das atividades econômicas.\item agricultura independe de clima e logística.\item o sul do estado é integralmente manguezal.\item grãos são produzidos apenas na ilha de São Luís.\item logística não interfere na competitividade agrícola.\end{enumerate}\end{questao}
\begin{questao}{\nivel{N4} 08.42 — Cebraspe (C/E)}A leitura integrada do território permite relacionar desmatamento, expansão agropecuária, disponibilidade hídrica e mudanças no uso da terra, temas que não se esgotam na descrição física do relevo.\end{questao}
\begin{questao}{\nivel{N4} 08.43 — FGV}Um município da Baixada apresenta cheias sazonais e ocupação urbana de áreas baixas. O planejamento deve priorizar:\begin{enumerate}[label=\Alph*)]\item ignorar hidrologia local.\item mapear risco, drenagem, ocupação e dinâmica sazonal das águas.\item tratar toda inundação como evento imprevisível.\item eliminar qualquer área úmida.\item usar apenas dados eleitorais.\end{enumerate}\end{questao}
\begin{questao}{\nivel{N4} 08.44 — Cebraspe (C/E)}O fato de o Porto do Itaqui ampliar capacidade de escoamento não implica que benefícios econômicos e infraestrutura social se distribuam automaticamente de forma uniforme no território.\end{questao}
\begin{questao}{\nivel{N4} 08.45 — FGV}A sequência “França Equinocial -- Beckman -- adesão à Independência -- Balaiada -- Vitorinismo” demonstra:\begin{enumerate}[label=\Alph*)]\item transição entre Colônia, Império e República.\item eventos todos do mesmo século.\item episódios exclusivamente militares.\item fatos sem relação com o Maranhão.\item apenas movimentos de independência.\end{enumerate}\end{questao}
\begin{questao}{\nivel{N4} 08.46 — Cebraspe (C/E)}A Batalha do Jenipapo pode ser cobrada em História do Maranhão porque o processo de Independência no norte do Brasil ultrapassou fronteiras provinciais atuais e afetou a adesão maranhense.\end{questao}
\begin{questao}{\nivel{N4} 08.47 — FGV}Ao estudar cultura maranhense em perspectiva geográfica, é correto relacioná-la:\begin{enumerate}[label=\Alph*)]\item à formação histórica, às matrizes populacionais e aos territórios de sociabilidade.\item apenas a entretenimento sem dimensão histórica.\item exclusivamente ao clima.\item apenas a indicadores econômicos.\item somente à legislação federal.\end{enumerate}\end{questao}
\begin{questao}{\nivel{N4} 08.48 — Cebraspe (C/E)}A distinção entre bacias limítrofes e bacias maranhenses ajuda a compreender fronteiras, drenagem interna e integração interestadual.\end{questao}
\begin{questao}{\nivel{N4} 08.49 — FGV}Um estudo sobre desenvolvimento sustentável no Maranhão deveria integrar, no mínimo:\begin{enumerate}[label=\Alph*)]\item ambiente, população, logística e estrutura produtiva.\item apenas o número de municípios.\item exclusivamente relevo.\item somente produção agrícola.\item apenas dados históricos do século XVII.\end{enumerate}\end{questao}
\begin{questao}{\nivel{N4} 08.50 — Cebraspe (C/E)}Compreender o Maranhão exige articular permanências históricas, desigualdades regionais, diversidade ambiental e inserção em cadeias logísticas nacionais e internacionais.\end{questao}


\section{Treino discursivo}
\begin{discursiva}[história, território e desenvolvimento]
Explique como a posição geográfica e a diversidade ambiental do Maranhão influenciaram sua ocupação econômica e sua integração logística, relacionando pelo menos dois elementos físicos e dois elementos econômicos.
\end{discursiva}

\section{Gabarito e resoluções comentadas — História e Geografia do Maranhão}

\begin{solucao}{08.01}\textbf{Certo.} A França Equinocial foi a tentativa francesa de colonização no Maranhão em 1612; França Antártica refere-se à experiência francesa no Rio de Janeiro no século XVI.\end{solucao}
\begin{solucao}{08.02}\textbf{Certo.} Guaxenduba, em 1614, integra a reação luso-portuguesa contra os franceses.\end{solucao}
\begin{solucao}{08.03}\textbf{Gabarito: B.} A ocupação holandesa de São Luís ocorreu em 1641, depois do ciclo francês do início do século XVII.\end{solucao}
\begin{solucao}{08.04}\textbf{Certo.} A Revolta de Beckman esteve ligada à crise do abastecimento, ao monopólio comercial e a conflitos envolvendo mão de obra indígena e jesuítas.\end{solucao}
\begin{solucao}{08.05}\textbf{Errado.} A adesão maranhense não foi imediata em 1822; consolidou-se em 1823 após conflitos regionais.\end{solucao}
\begin{solucao}{08.06}\textbf{Gabarito: B.} A Balaiada foi revolta regencial de forte base popular, com causas sociais e políticas.\end{solucao}
\begin{solucao}{08.07}\textbf{Certo.} Esses são os limites interestaduais e marítimo do Maranhão.\end{solucao}
\begin{solucao}{08.08}\textbf{Certo.} A distinção aparece expressamente no conteúdo editalício.\end{solucao}
\begin{solucao}{08.09}\textbf{Gabarito: B.} APA é unidade de uso sustentável; parque nacional é unidade de proteção integral.\end{solucao}
\begin{solucao}{08.10}\textbf{Certo.} Babaçu é elemento marcante das formações de cocais.\end{solucao}

\begin{solucao}{08.11}\textbf{Gabarito: A.} La Touche liderou a expedição francesa ligada à França Equinocial e à fundação de São Luís.\end{solucao}
\begin{solucao}{08.12}\textbf{Certo.} Jenipapo ocorreu no Piauí, mas integra o processo regional de consolidação da Independência no norte.\end{solucao}
\begin{solucao}{08.13}\textbf{Certo.} Beckman contestava mecanismos do sistema colonial, não constituía projeto nacional de independência.\end{solucao}
\begin{solucao}{08.14}\textbf{Gabarito: B.} O Vitorinismo deriva da liderança e influência política de Vitorino Freire no século XX.\end{solucao}
\begin{solucao}{08.15}\textbf{Certo.} O episódio ocorreu em contexto de intensa disputa política e mobilização social.\end{solucao}
\begin{solucao}{08.16}\textbf{Gabarito: C.} O Maranhão é espaço de transição ambiental e territorial, o que explica sua diversidade de formações.\end{solucao}
\begin{solucao}{08.17}\textbf{Certo.} Regime de chuvas condiciona rios, vegetação e atividades agropecuárias.\end{solucao}
\begin{solucao}{08.18}\textbf{Certo.} Hidrografia é conteúdo funcional: abastecimento, ocupação, produção e ambiente.\end{solucao}
\begin{solucao}{08.19}\textbf{Gabarito: A.} A Baixada é marcada por planícies e áreas sujeitas a inundações sazonais.\end{solucao}
\begin{solucao}{08.20}\textbf{Certo.} O sul e leste do Maranhão integram a expansão agrícola do MATOPIBA.\end{solucao}
\begin{solucao}{08.21}\textbf{Gabarito: B.} O complexo portuário e as ferrovias são elementos centrais dos corredores de exportação.\end{solucao}
\begin{solucao}{08.22}\textbf{Certo.} São Luís concentra serviços superiores, administração, porto e atividades industriais.\end{solucao}
\begin{solucao}{08.23}\textbf{Certo.} Imperatriz exerce forte centralidade regional no sudoeste maranhense.\end{solucao}
\begin{solucao}{08.24}\textbf{Gabarito: A.} Babaçu = extrativismo vegetal; bumba-meu-boi = cultura; Itaqui = logística.\end{solucao}
\begin{solucao}{08.25}\textbf{Certo.} Urbanização pode coexistir com desigualdade territorial de infraestrutura e serviços.\end{solucao}

\begin{solucao}{08.26}\textbf{Certo.} Guaxenduba pertence ao conflito contra franceses, não à expulsão holandesa.\end{solucao}
\begin{solucao}{08.27}\textbf{Gabarito: B.} A ordem correta é França Equinocial, Guaxenduba, ocupação holandesa e Beckman.\end{solucao}
\begin{solucao}{08.28}\textbf{Certo.} A Balaiada combinou disputas políticas e profundas tensões sociais, com ampla participação popular.\end{solucao}
\begin{solucao}{08.29}\textbf{Certo.} A Revolução de 1930 reorganizou a relação entre centro e estados e alterou os grupos políticos locais.\end{solucao}
\begin{solucao}{08.30}\textbf{Gabarito: A.} Os eventos atravessam disputa colonial, crise do sistema colonial e Regência.\end{solucao}
\begin{solucao}{08.31}\textbf{Certo.} A localização de transição ajuda a explicar gradientes de chuva e mosaico vegetal.\end{solucao}
\begin{solucao}{08.32}\textbf{Gabarito: A.} Babaçuais são característicos da Mata dos Cocais.\end{solucao}
\begin{solucao}{08.33}\textbf{Certo.} Manguezais são ecossistemas costeiro-estuarinos.\end{solucao}
\begin{solucao}{08.34}\textbf{Gabarito: B.} Essa é a distinção central entre as duas categorias de unidade de conservação.\end{solucao}
\begin{solucao}{08.35}\textbf{Certo.} Planejamento territorial deve integrar fatores físicos e riscos ambientais.\end{solucao}
\begin{solucao}{08.36}\textbf{Gabarito: A.} A EFC integra o corredor de Carajás até a região portuária de São Luís.\end{solucao}
\begin{solucao}{08.37}\textbf{Certo.} A estrutura econômica é diversificada e territorialmente desigual.\end{solucao}
\begin{solucao}{08.38}\textbf{Certo.} Crescimento econômico não implica automaticamente redução de desigualdades.\end{solucao}
\begin{solucao}{08.39}\textbf{Gabarito: A.} São Luís e Imperatriz são centros regionais com funções diferentes e áreas de influência próprias.\end{solucao}
\begin{solucao}{08.40}\textbf{Certo.} As manifestações culturais resultam de matrizes africanas, indígenas e europeias articuladas historicamente.\end{solucao}

\begin{solucao}{08.41}\textbf{Gabarito: A.} Atividade agrícola depende de condições naturais, tecnologia, terra, crédito e acesso logístico aos mercados.\end{solucao}
\begin{solucao}{08.42}\textbf{Certo.} Uso da terra é tema integrado de geografia física, econômica e ambiental.\end{solucao}
\begin{solucao}{08.43}\textbf{Gabarito: B.} Planejamento deve considerar dinâmica de cheias, ocupação e drenagem, reduzindo exposição e vulnerabilidade.\end{solucao}
\begin{solucao}{08.44}\textbf{Certo.} Infraestrutura logística pode gerar ganhos concentrados e não assegura distribuição territorial uniforme.\end{solucao}
\begin{solucao}{08.45}\textbf{Gabarito: A.} A sequência atravessa Colônia, Império e República.\end{solucao}
\begin{solucao}{08.46}\textbf{Certo.} A consolidação da Independência ocorreu em escala regional e envolveu províncias vizinhas.\end{solucao}
\begin{solucao}{08.47}\textbf{Gabarito: A.} Cultura é expressão espacial e histórica das relações sociais e identidades coletivas.\end{solucao}
\begin{solucao}{08.48}\textbf{Certo.} A distinção permite ler tanto limites estaduais quanto drenagem interna.\end{solucao}
\begin{solucao}{08.49}\textbf{Gabarito: A.} Desenvolvimento sustentável exige visão multidimensional do território.\end{solucao}
\begin{solucao}{08.50}\textbf{Certo.} Essa é a síntese adequada para interpretar o estado em perspectiva histórica e geográfica integrada.\end{solucao}

